\scanpage{127}
\begin{center}- 145 -\end{center}
\margin{n\textsuperscript{o} 316}

Si $x\longrightarrow x_M$ est une représentation irréductible de dimension finie, si $\beta$ est la forme bilinéaire associée, et si $\dim\unsure{h|\mathfrak h}$ n'est pas divisible par la caractéristique de $K$, montrer que $H^p(\unsure{g|\mathfrak g},M)=\{0\}$ pour tout $p$. (Utilisant la prop.12, montrer que $\unsure{\mathbf T|\mathrm T}$ est un automorphisme de $C^*(\unsure{g|\mathfrak g},M)$, donc induit des automorphismes de $Z^p(\unsure{g|\mathfrak g},M)$ et $B^p(\unsure{g|\mathfrak g},M)$. Pour $f\in Z^p(\unsure{g|\mathfrak g},M)$, $\begin{aligned}\unsure{\mathbf T|\mathrm T}f&=\\&=(d\rho+\rho d)f=d\rho f\in B^p(\unsure{g|\mathfrak g},M)\end{aligned}$, d'où $Z^p(\unsure{g|\mathfrak g},M)=B^p(\unsure{g|\mathfrak g},M)$).
\figure{\begin{tikzpicture}\draw[line width=.6pt] (0,0)--(2.5,0);\end{tikzpicture}}
\margin{126}
